% !TEX root = ../main.tex
% ESEMPIO 6: slide del titolo con più autori, loghi e ringraziamenti ai finanziamenti.
%   - \begingroup ... \endgroup: i metadati (\title, \author, ...) valgono solo qui, così il titolo vero della
%     talk, scritto in main.tex, non cambia. Nella tua presentazione questo blocco si scrive in main.tex.
%   - Sotto \titlepage: i loghi in una tabella dentro adjustbox (altezza al massimo il 12,5% della slide) e,
%     in piccolo, il testo dei finanziamenti con \parMini (definita in 00_macro.tex).
%   - Qui al posto dei loghi ci sono due immagini di esempio: sostituiscile con le tue.

\begingroup
\title{Example: title slide with logos and funding}
\author{\small \textbf{Speaker}\inst{1} \and Author Two\inst{1} \and Author Three\inst{2}}
\institute{\inst{1} Affiliation 1 \and \inst{2} Affiliation 2}
\date{}

\begin{frame}
    \titlepage

    \vfill

    \vspace{-2em}
    \begin{adjustbox}{max totalsize={\textwidth}{.125\textheight},center}
        \begin{tabular}{lr}
            \includegraphics[width=.05\textwidth]{assets/img/patrick_happy.jpg} &
            \includegraphics[width=.1\textwidth]{assets/img/patrick_scientist.jpg}
        \end{tabular}
    \end{adjustbox}

    \parMini{The work was partially supported by
        \emph{i)} project A, funded by the funder under the call number 1;
        \emph{ii)} project B, funded by another funder.
    }

\end{frame}
\endgroup
